% Source: https://artofproblemsolving.com/community/c3h1600336_greatest_number_smaller_than_1_but_not_1
% Stable identity: aops-topic-1600336-q01
Let $x<1$ and put $y=\dfrac{x+1}{2}$, the midpoint of $x$ and $1$. Then
\[
y-x=\dfrac{1-x}{2}>0\qquad\text{and}\qquad 1-y=\dfrac{1-x}{2}>0,
\]
so $x<y<1$. Hence no $x<1$ can be the greatest real number below $1$: $y$ is larger and still below $1$. $\blacksquare$

\textbf{Remarks.} In the same way there is no greatest rational number below $1$ (the midpoint of two rationals is rational). The number $0.999\ldots$ is not a candidate either: it is the limit of $0.9,0.99,0.999,\dots$, and that limit equals $1$, because $1-0.99\ldots9$ (with $n$ nines) $=10^{-n}$ becomes smaller than any positive number. Other number systems (for example the hyperreal numbers, which contain infinitesimals) still have no greatest element below $1$, by the same midpoint argument.
